\section{Validación, pruebas y observabilidad}

\subsection{Invariantes por dominio}

Los reductores deben proteger invariantes baratas y deterministas. Las reglas que necesitan acceso a escena, red o base de datos pertenecen a servicios previos al despacho.

\begin{table}[H]
\centering
\caption{Invariantes mínimas recomendadas}
\begin{tabularx}{\textwidth}{>{\bfseries}p{3cm}X X}
\toprule
Dominio & Invariante & Comportamiento esperado \\
\midrule
Player & Cada dimensión del collider es al menos 0.01. & Limitar en el reductor; ya implementado. \\
Player & El nivel no es negativo. & Limitar o rechazar la acción según la regla de negocio. \\
Player & La posición contiene valores finitos. & Rechazar NaN e infinito antes de actualizar. \\
World & Ancho y alto son positivos. & Rechazar mundos sin área válida. \\
Decoration & El ID no está vacío y no se duplica. & Validar al agregar la instancia. \\
Decoration & Inventario solo se usa cuando está habilitado. & Vaciar, ignorar o conservar según una política documentada. \\
Todos & IDs y textos nulos se normalizan. & Convertir a cadena vacía o rechazar según obligatoriedad. \\
\bottomrule
\end{tabularx}
\end{table}

\subsection{Pirámide de pruebas}

\subsubsection{Pruebas unitarias de reductores}

Son la base del sistema. Cada tipo de acción debe probar estado inicial, payload, resultado y preservación de campos no relacionados. También deben cubrir valores límite y acciones desconocidas.

\subsubsection{Pruebas de stores}

Verifican que \code{StateChanged} se emite una sola vez cuando hay cambio y no se emite para una transición equivalente. Deben comprobar que el evento entrega la acción que causó la transición.

\subsubsection{Pruebas de integración en Play Mode}

Cubren la conexión con \code{Rigidbody2D}, colliders, \code{WorldYSort}, animación y cambios de escena. Casos esenciales:

\begin{itemize}
  \item el jugador no atraviesa obstáculos y puede deslizarse por sus bordes;
  \item la reconciliación conserva alineados store y cuerpo físico;
  \item cambiar el tamaño del collider actualiza el componente y respeta el mínimo;
  \item cambiar de escena preserva la selección de rol;
  \item crear o quitar una decoración sincroniza el registro visual sin duplicados.
\end{itemize}

\subsection{Registro de transiciones}

Durante desarrollo, un observador opcional puede registrar \code{timestamp}, store, tipo de acción y resumen del estado resultante. No debe incluir correo, tokens ni inventarios sensibles. Para estados grandes se registra un resumen o diferencia, no la instantánea completa en cada frame.

\subsection{Rendimiento}

El patrón no implica despachar todos los valores visuales continuamente. Recomendaciones:

\begin{itemize}
  \item no guardar en el estado global cuadros de animación que cambian cada frame, salvo que sean necesarios para red o repetición;
  \item evitar reconstruir colecciones completas para una sola decoración en escenarios masivos; usar actualización por ID y estructuras adecuadas;
  \item no publicar eventos si \code{StatesAreEqual} confirma que no hubo cambio;
  \item mantener la física en \code{FixedUpdate} y la lectura de entrada en \code{Update};
  \item perfilar antes de optimizar o reemplazar estructuras claras por soluciones complejas.
\end{itemize}

\subsection{Criterios de aceptación de la arquitectura}

La implementación objetivo se considera completa cuando:

\begin{enumerate}
  \item todo estado se modifica únicamente mediante acciones y reductores;
  \item los estados planeados tienen tipos C\# concretos y serializables;
  \item cada campo tiene propietario, significado, unidad y valor inicial documentados;
  \item los efectos de Unity están fuera de los reductores;
  \item las transiciones críticas cuentan con pruebas unitarias y de integración;
  \item el guardado y la carga reconstruyen una instantánea equivalente;
  \item los IDs resuelven assets y objetos sin serializar referencias de escena.
\end{enumerate}
